%========================================================== TECHNICAL GAPS
\section{Technical gap register}\label{sec:gaps}

\subsection{First, what Betaflight actually is}

Several gaps below are ``Betaflight things'', so here is the minimum needed to read them.

\begin{explain}[The flight controller, in plain terms]
The \textbf{flight controller} is a small circuit board on the drone with a gyroscope on it, running
firmware called \textbf{Betaflight}. Its job is narrow and fast: take four numbers --- throttle,
roll, pitch, yaw --- several thousand times a second, and spin the four motors so the drone does what
those numbers ask. It handles the violent, millisecond-scale balancing that keeps a quadcopter in
the air. We do not have to write that, and we could not do it better.

Normally those four numbers come from a \textbf{radio receiver} wired to the flight controller,
listening to a human pilot's remote control.

\textbf{MSP} (MultiWii Serial Protocol) is simply a message format for talking to the flight
controller down a serial wire. Betaflight speaks it. The feature we need is called \textbf{MSP
override}: while a switch on the pilot's remote is flipped on, Betaflight takes its four stick
numbers from MSP messages --- from our Jetson --- instead of from the radio. Flip the switch off and
the pilot has the drone back, immediately, in firmware. That is our kill path.

Two more terms that appear below:
\begin{itemize}
\item \textbf{ACRO and ANGLE} are flight modes. In \textbf{ACRO}, stick position means \emph{rate of
rotation}: hold the stick over and the drone keeps rolling; centre it and it stops rolling but stays
tilted wherever it ended up. In \textbf{ANGLE}, stick position means \emph{lean angle}: centre the
stick and it returns to level by itself. Our trained policy outputs rotation rates, so it needs
ACRO. Plan~B wants ANGLE, because ``hold still and hover'' is then almost free.
\item \textbf{Blackbox} is Betaflight's built-in flight recorder. It logs throttle, gyro, motor
outputs and battery voltage during a flight, and you download it over USB afterwards. Almost every
measurement we need about the drone comes out of it.
\end{itemize}
\end{explain}

\subsection{Gap register}

Effort is one person's working time. ``Blocks'' names what cannot start until this is done.

{\setlength{\tabcolsep}{3pt}\begin{longtable}{@{}p{0.85cm} p{3.9cm} p{6.5cm} p{1.45cm} p{1.95cm}@{}}
\toprule
ID & Gap & What it means, and why it matters & Effort & Blocks \\
\midrule\endhead
\multicolumn{5}{@{}l}{\textbf{A --- Flight controller and the link to it} \emph{(much reduced: the organizers ship the protocol layer)}}\\
\midrule
A0 & Copy \code{~/target/} off a board & \code{scp -r dcl@192.168.55.1:~/target ./}. Until somebody does this, nobody else can develop against the fake flight controller. & 10 min & All adapter work \\
A1 & Adapter, not protocol & MSP framing is provided. We write the layer that turns policy output into \code{RCTransmitter.set\_control()} calls. & 4\,h & Everything autonomous \\
A2 & Free the serial port & \code{sudo setup\_jetson\_uart.sh -{}-apply}, once per board. Linux runs a login console on that port by default. \textbf{Skipping this looks exactly like a dead flight controller.} & 5 min & Any link at all \\
A3 & Rate curve inversion & \textbf{Still ours.} Betaflight maps stick to rotation rate along a curve (\code{rc\_rate}, \code{super\_rate}, \code{expo}), gentle in the middle and aggressive at the ends. Our policy speaks radians per second. A linear approximation is wrong by tens of percent mid-range. & 4\,h & Accurate control \\
A4 & Throttle curve inversion & \textbf{Still ours.} Same for throttle (\code{thr\_mid}, \code{thr\_expo}). & 1\,h & Height control \\
A5 & Sign and axis mapping & \textbf{Still ours, and the most dangerous item here.} Isaac is forward-left-up, Betaflight is forward-right-down, plus channel order. A reversed sign flips the drone on takeoff. Verify one axis at a time, props off. & 2\,h & Any flight \\
A6 & Watchdog: configure and verify & \code{RCTransmitter} already has one, in its own thread. Set the timeout, then prove it against \code{fake\_fc.py}'s frozen mode. & 1\,h & Safety sign-off \\
A7 & \textbf{Loop rate mismatch} & Organizers: attitude at 30--50\,Hz is realistic, a hard real-time loop is not. Our policy runs at 60\,Hz and eats gyro as input. Measure, then match the simulator to it. & 2\,h & \textbf{Training validity} \\
A8 & Gyro units unverified & Sources conflict: raw counts at 16.4 per \si{\degree\per\second} versus already in \si{\degree\per\second}. \textbf{Verify empirically} --- the wrong choice is a silent 16$\times$ error. & 30 min & Sensor fusion \\
A9 & Arming flow & \code{msp\_bench.py demo} clears an MSP arming lock, arms, then re-asserts it; Betaflight also blocks arming while an MSP client is connected. Read their code and copy the sequence. & 2\,h & Any real run \\
A10 & Override mask: verify, do not assume & The mask\,=\,11 finding (throttle not overridden) came from \emph{another team's} dump. Since \code{demo} drives throttle successfully, ours may already be right. Check \code{diff all} on our drone. & 10 min & Autonomous hover \\
A11 & ANGLE mode not assigned & Flies ACRO; Plan~B needs ANGLE. CLI-assignable --- confirm it is permitted. & 15 min & Plan~B \\
A12 & Blackbox not configured & Must be on and recording gyro debug data, or the measurement flight yields nothing. & 15 min & All diagnostics \\
A13 & Camera/IMU time sync & \textbf{No hardware sync line exists on this carrier.} Must be solved in software, and the exposure must be pinned first or auto-exposure moves the capture instant every frame. & 4\,h & Latency measurement \\
\midrule
\multicolumn{5}{@{}l}{\textbf{B --- Onboard software}}\\
\midrule
B1 & Capture with \emph{real} timestamps & Viewers exist. But \code{cv2.VideoCapture} discards the hardware capture timestamp and gives arrival time with tens of ms of jitter. Use the GStreamer path where timing matters. & 4\,h & Perception \\
B1b & \textbf{Image processor needs a display} & Colour, auto-exposure and white balance run in NVIDIA's ISP, which needs a graphics context a plain SSH session lacks. Headless gives flat, un-exposed frames \emph{by design}. Arrange a display context, or design for minimally-processed frames. & 4\,h & Usable images \\
B2 & \textbf{No framework on the board} & Confirmed: OpenCV, NumPy, pyserial, python3-gi --- no PyTorch. The NumPy policy runtime is now \emph{mandatory}. The detector must go to ONNX via OpenCV's DNN module or TensorRT. Never \code{pip install opencv-python}: it silently breaks the camera pipeline. & 3\,h & Plan~A \\
B3 & No gate counter & The ``which gate is next'' input has no source on the real drone. & 1 day & Plan~A \\
B4 & No takeoff routine & The policy was trained starting \emph{in the air} near a gate. The real race starts on the floor. Something must take off, stabilise, and hand over. & 4\,h & Any real run \\
B5 & Crash and abort logic is simulator-era & The existing logic auto-re-arms after a crash, which is a simulator behaviour and dangerous on a real drone. & 3\,h & Safety sign-off \\
B6 & No onboard logging standard & Every flight needs a synchronised record of inputs, outputs and detections, or we cannot diagnose anything. & 2\,h & All learning \\
B7 & Compute budget unverified & Detector plus policy plus capture, at 25\,W, at the target frame rate. Unknown until measured. & 2\,h & Loop rate \\
\midrule
\multicolumn{5}{@{}l}{\textbf{C --- Perception}}\\
\midrule
C1 & Detector has never seen a real gate & Existing detectors were tuned on simulator images. See Section~\ref{sec:detect}. & 2 days & Any autonomy \\
C2 & Camera not calibrated & The spec states no calibration will be provided, and it is a screw-mount lens. & 1\,h & Perception accuracy \\
C3 & Camera tilt unknown & It is adjustable, even in flight, so it is a choice we must make and then measure. A \SI{5}{\degree} error took crashes to 32 per 100 gates. & 30 min & Perception accuracy \\
C4 & \textbf{Sensor mode may not exist} & The carrier wires only \textbf{two of the camera's four data lanes}, capping resolution $\times$ framerate. \code{v4l2-ctl -{}-list-formats-ext} is the truth, not the datasheet. The \SI{75}{\degree} figure is quoted for 1920$\times$1080 at 60\,fps --- if that mode is missing, our field of view is wrong before we start. & 10 min & Everything perception \\
\midrule
\multicolumn{5}{@{}l}{\textbf{D --- Simulator fidelity}}\\
\midrule
D1 & Thrust-to-weight is a guess & Hover throttle assumed 25.5\,\%. Worth 58 crashes per 100 gates if wrong by 10\,\%. & 5 min & Training validity \\
D2 & Inertia and rate response wrong & 5-inch airframe values on an 8-inch drone; see the warning in the status section. & 4\,h & Training validity \\
D3 & Zero latency modelled & The most damaging single factor found: \SI{50}{ms} alone took crashes from 2 to 55 per 100 gates. & 3\,h & Training validity \\
D4 & No drag model & Significant above \SI{10}{m/s}. & 3\,h & Training validity \\
D5 & Perfect vision assumed & The simulator hands the policy exact gate corners with no noise, no dropouts and never the wrong gate. This is the largest single difference between simulation and reality. & 1 day & Transfer \\
D6 & Course map: official table in hand & The official coordinates confirm $85\times165$\,ft and our corrected track matches exactly. \textbf{Still survey the as-built placement} --- a 7.5\,\% scale error alone took crashes from 2 to 15.5 per 100 gates. & 2\,h & Training validity \\
D7 & \textbf{Policy rate may be too fast} & Simulator runs a 60\,Hz policy; real gyro feedback may be 30--50\,Hz. Set decimation to the measured rate \emph{before} the long training run. & 1\,h & Training validity \\
\midrule
\multicolumn{5}{@{}l}{\textbf{E --- Operations}}\\
\midrule
E1 & No compiler on the laptop & \textbf{Downgraded.} The Betaflight software-in-the-loop build is no longer needed --- the organizers' \code{fake\_fc.py} is pure Python and does the job better. & 10 min & nothing now \\
E2 & Practice slots unknown & Flight time is rationed and scheduled. Everything must be planned around the slot list. & --- & All planning \\
E3 & Battery and propeller count unknown & The real limit on how many attempts we get. & --- & Risk budget \\
\bottomrule
\caption{Technical gap register as of 17 September, after reading the organizers' Orin NX guides. Nothing here is a criticism of the existing code --- almost all of it was written for the virtual qualifier, which was a different race with a different interface. The A-series shrank considerably once it became clear the organizers ship the protocol layer themselves.}\label{tab:gaps}
\end{longtable}}

%========================================================== WORKSTREAMS
\section{The plan, organised as parallel workstreams}\label{sec:ws}

The mistake to avoid is queueing everything behind the flight-controller link. Only one workstream
actually depends on it. Three others need no software at all and should start the same hour.

\begin{table}[H]\centering\footnotesize
\begin{tabularx}{\linewidth}{@{}p{1.1cm} p{3.0cm} X p{2.4cm}@{}}
\toprule
Stream & Name & Content & Waits for \\
\midrule
\textbf{WS-1} & Link and control & MSP implementation, curve inversion, watchdog, bench validation (gaps A1--A12) & \textbf{Nothing.} This is the critical path \\
\textbf{WS-2} & Vehicle characterisation & Weigh, balance, measure; the piloted measurement flight; hover thrust, rate response, drag, battery sag (D1--D4) & \textbf{Nothing} --- a human pilot and a blackbox log need no code \\
\textbf{WS-3} & Perception & Camera bring-up, calibration, data capture, labelling, detector, gate counter (B1, B3, C1--C4, D5) & Camera only \\
\textbf{WS-4} & Course survey & Tape-measure every gate from a fixed datum; reconcile with both documents; build the map (D6) & \textbf{Nothing} --- hall access and a tape measure \\
\textbf{WS-5} & Simulator and training & Fold measurements in; cloud training with randomised unknowns; evaluate with the stress harness & \textbf{Nothing} --- starts with wide randomisation \\
\textbf{WS-6} & Integration and flight test & The bring-up ladder, in practice slots & WS-1, plus WS-3 and WS-4 for gate work \\
\textbf{WS-7} & Plan~B & The slow stop-and-line-up controller & WS-1, WS-3 \\
\textbf{WS-0} & Operations & Slots, batteries, organizer questions, installs, logistics & \textbf{Nothing} \\
\bottomrule
\end{tabularx}
\caption{Eight workstreams. Five of the eight can start immediately and in parallel.}
\end{table}

\begin{takeaway}[How to staff this]
Put your strongest systems person on \textbf{WS-1} alone, with a second person reviewing the sign
and axis conventions --- a reversed sign there flips the drone on takeoff, and it is the single most
common first-flight failure in this whole field.

\textbf{WS-2} needs the pilot plus one engineer, and it is the cheapest high-value work available:
one battery, one log, and the simulator stops guessing.

\textbf{WS-4} is two people and ninety minutes, and it retires a risk worth 15 crashes per 100
gates. Do it on day one and do it again before the scored days.

\textbf{WS-3} is the longest pole after WS-1 and it is mostly patient work --- capture, label, train,
measure. Start the capture the hour the camera produces a frame.
\end{takeaway}

%========================================================== PROCEDURES
\section{Field procedures}\label{sec:proc}

\subsection{P1 --- Camera calibration}
\begin{enumerate}
\item Confirm the sensor mode first (gap C4). The \SI{75}{\degree} figure applies to 1920$\times$1080
at 60\,fps. If our pipeline captures another mode, the field of view is different before we start.
\item Print a checkerboard or ChArUco target and \textbf{glue it to something rigid and flat}. A
curled printout produces a confident, wrong calibration.
\item Capture 30--50 views at the race exposure and gain settings: target near and far, tilted in
both axes, in all four corners of the frame.
\item Solve for focal length, principal point and distortion (OpenCV). Expect a reprojection error
below about half a pixel; if it is worse, the target moved or the board is not flat.
\item \textbf{Sanity check independently:} tape measure flat against a wall, camera exactly
\SI{2.00}{m} away, and confirm the width spanned matches the focal length you solved for.
\item Record the result with the date and the specific drone. The four airframes may not match.
\end{enumerate}

\subsection{P2 --- Camera tilt}
\begin{enumerate}
\item Set the drone on a surface you have confirmed is level with a spirit level.
\item Measure the camera mount angle with a phone inclinometer held against the lens face or the
bracket. Repeat three times and take the median; this is easy to get wrong by a couple of degrees.
\item Cross-check optically: place a target at a known height at a known distance and confirm it
lands where the calibrated model predicts.
\item Tilt is adjustable, so \textbf{decide} it: a lower tilt (around \SI{10}{\degree}) suits Plan~B's
slow, near-level flight; a higher tilt (around \SI{20}{\degree}) suits fast forward-leaning flight.
Choose, set, measure, and write it down. Then use that number in the simulator.
\end{enumerate}

\subsection{P3 --- Vehicle diagnostics: one battery, one log}\label{sec:diag}

This is the highest-value thirty minutes of the week. Fly it with a human pilot, with blackbox
logging switched on, in this order. It needs no autonomy software at all.

\begin{table}[H]\centering\footnotesize
\begin{tabularx}{\linewidth}{@{}p{0.6cm} p{4.4cm} X@{}}
\toprule
\# & Manoeuvre & What it yields \\
\midrule
1 & Armed on the ground, 10\,s idle & Gyro noise floor; vibration spectrum; motor idle value \\
2 & Stable hover at \SI{1.5}{m} for 20\,s & \textbf{Hover throttle} --- gap D1, the most valuable number available. Also a ground-effect reference \\
3 & Three sharp rolls, three pitches, three yaws, settling between each & \textbf{Rate response}: rise time, overshoot, peak angular acceleration, per axis --- gap D2. This is what replaces guessing the inertia \\
4 & Full-throttle vertical climb, 1.5\,s, then recover & \textbf{Maximum thrust} and true thrust-to-weight \\
5 & Hold about \SI{20}{\degree} of pitch down a straight until the speed stops rising & \textbf{Drag} --- gap D4. Drag coefficient $\approx m g \tan\theta / v^2$ \\
6 & Hover again for 20\,s near the end of the pack & \textbf{Battery sag}: how much hover throttle rises as the pack empties \\
7 & Land, disarm, download the log, and save a \code{diff all} from the CLI & The settings dump. Do not skip this; it is the thing that will be missing at 2\,a.m.\ \\
\bottomrule
\end{tabularx}
\end{table}

Repeat manoeuvre~2 on every airframe you might fly. Hover throttle is the number the whole simulator
hangs from, and four supposedly identical drones are rarely identical.

\subsection{P4 --- Latency}
Latency is the most damaging single factor in the durability study, and it is also the one people
estimate instead of measuring. Measure it.
\begin{enumerate}
\item Synchronise the clocks first (gap A12), or the answer is meaningless.
\item Put a bright LED in the camera's view, wired so the Jetson can switch it.
\item Flash it, and record: the time the Jetson commanded the flash; the timestamp of the first frame
showing it lit; the time the detector produced an output; the time the command went out over MSP; and
the time the gyro first responded.
\item That chain gives you each stage separately, which matters --- knowing the total is \SI{90}{ms}
is useful, knowing that \SI{60}{ms} of it is the detector tells you what to fix.
\item Feed the measured total, and its variation, into the simulator as a randomised band.
\end{enumerate}

\subsection{P5 --- Course survey}
\begin{enumerate}
\item Choose a datum: a permanent corner as the origin, plus a second fixed feature to define the
direction of the $x$ axis. Mark both with tape and photograph them.
\item For each gate, measure the \textbf{base of both uprights} from the datum. Two points give you
the centre \emph{and} the heading in one operation, which is far more reliable than measuring an
angle directly.
\item Record the bottom-bar height for one gate of each type, and confirm the double gate's upper
opening centre.
\item Note what the drawing does not show: pylons, netting, lighting rigs, anything above \SI{4}{m}.
\item Convert to the simulator's frame and difference against \emph{both} the organizer table and the
drawing. Reconcile every gate that differs by more than \SI{0.25}{m} before training on it.
\item Re-measure every morning and before every scored run. Ask what tolerance gates are re-placed
to; that number becomes the gate-jitter range in training.
\end{enumerate}

\subsection{P6 --- Gate detection and data labelling}\label{sec:detect}

This deserves its own treatment, because it is the largest difference between our simulator and
reality, and it is the workstream most likely to be underestimated.

\begin{explain}[Why this is hard, when the simulator makes it look easy]
In the simulator, the policy is handed the eight corners of the gate as exact pixel coordinates,
computed by projecting the gate's known geometry through a perfect pinhole camera. There is no noise.
Corners are never missed. It is never the wrong gate. Visibility is known exactly.

On the real drone, every one of those is false. You have a compressed image of an orange-and-white
frame against a cluttered hall, with motion blur, a rolling shutter that skews fast-moving objects,
changing light, other gates in frame, orange pylons that look like gate material, and a detector that
will sometimes return nothing and occasionally return something confidently wrong.

\textbf{Closing this gap has two halves, and the second is the one teams skip:} make the detector as
good as is reasonable, \emph{and} make the policy tolerant of a realistic detector by reproducing the
detector's \emph{measured} error statistics inside the simulator. Perfect detection is not
achievable; a policy that degrades gracefully when detection is imperfect is.
\end{explain}

\paragraph{Stage 0 --- classical detector, first hour.}
Tune the existing colour-and-geometry detector on real frames --- it lives in \code{vision/} as
\code{snake\_gate\_detector}, and \code{tools/hsv\_tuner} adjusts its colour thresholds. No data, no
training, no installation. It will be noisier
and shorter-ranged than a learned detector, but it is available immediately and it is the fallback
that cannot fail to be ready. Plan~B tolerates a noisy detector far better than Plan~A does.

\paragraph{Stage 1 --- capture, starting the hour the camera works.}
Data capture is the bottleneck, so begin it before you need it. Capture must cover:
\begin{itemize}
\item distances from about \SI{1}{m} to \SI{15}{m}, and all approach angles including oblique;
\item \emph{every} gate on the course --- the backgrounds differ, and that is what a detector learns;
\item gates partially cut off by the frame edge, and two gates overlapping;
\item the double gate, from below and from the upper approach;
\item backlit and shadowed conditions, and whatever the lighting does across a day;
\item \textbf{negatives}: frames with no gate, and frames containing confusers --- orange pylons, red
signage, other drones, people;
\item \textbf{motion blur}, which means recording every manual pilot flight. Hand-carried footage is
easy to collect but is missing exactly the degradation that matters in flight.
\end{itemize}

\paragraph{Stage 2 --- labelling, done cheaply.}
Labelling is where teams lose a day. Three techniques cut the work by roughly five to ten times:
\begin{enumerate}
\item \textbf{Auto-label, then correct.} Run the classical detector over all footage and keep its
output as draft labels. Correcting a drawn box is far faster than drawing one.
\item \textbf{Geometric auto-labelling --- the big win.} We know the gate's exact geometry
(\SI{2700}{mm} outer, \SI{1500}{mm} inner) and, after the survey, exactly where each gate is. If the
camera pose for a frame is known, the eight corners project into the image exactly --- the label is
free and pixel-perfect. Bootstrap it: hand-label around 40 frames, solve the camera pose from those
by PnP, then track along the video and re-project labels for every frame in between. This is how an
afternoon produces thousands of labelled frames instead of a few hundred.
\item \textbf{Label what the policy actually needs.} The observation wants eight keypoints plus
visibility flags --- not bounding boxes. Keypoint labelling is slow. A hybrid is cheaper: let a
learned model find the gate as a box, then fit the quadrilateral inside it with classical line
fitting. Much less labelling, and \code{vision/yolo\_hybrid\_gate\_detector.py} is already shaped
this way.
\end{enumerate}
For a single, visually distinctive object class, fine-tuning a pretrained detector typically needs
\textbf{300--800 well-chosen labelled frames}, not tens of thousands. Spend the effort on coverage
and on the hard cases, not on volume.

\paragraph{Stage 3 --- train and deploy.}
Fine-tune on the laptop GPU; it is a small model and takes minutes. Then build the accelerated engine
\textbf{on the Jetson itself} --- these are specific to the device and software version and cannot be
prepared in advance. Budget real time for this at the venue.

\paragraph{Stage 4 --- measure the detector, because this is what feeds the simulator.}
The useful metrics are not the ones detector papers report. Measure:
\begin{itemize}
\item detection rate against distance --- what \emph{is} our reliable range?
\item corner pixel error against distance and viewing angle, as a median and a 95th percentile;
\item false positives per frame, and specifically what triggers them;
\item how often the identity is wrong;
\item throughput and latency at \SI{25}{W}.
\end{itemize}
Get ground truth by placing the drone at tape-measured positions relative to a gate: those frames have
exactly known corner positions, so the error can be measured rather than estimated.
\code{tools/eval\_detectors.py}, \code{tools/eval\_gate\_pose.py} and \code{tools/vision\_rate.py}
exist for this.

\paragraph{Stage 5 --- close the loop.}
Feed those measurements back as the simulator's vision model: pixel noise as a function of distance,
dropout rate, maximum range, frame rate, latency. Then retrain. This is the step that converts a
detector measurement into a policy that survives a real detector, and it is the reason Stage~4 is not
optional.

\paragraph{The gate counter.}
Separately from detection: maintain the expected gate number, predict where that gate should appear
using the surveyed map, and accept a detection only if it lands near the prediction. Advance the
count on a confirmed pass, with a short lockout so one crossing cannot count twice. \textbf{Encode
explicitly that the double gate counts twice per lap} --- one physical gate carrying two different
numbers is where a counter desynchronises. If the expected gate is not seen for a while, widen the
search, then fall back to a hold; never silently re-label, and log every re-association. Test all of
this offline by replaying a hand-carried lap before it ever flies.

%========================================================== SCHEDULE
\section{Schedule and decision gates}\label{sec:schedule}

\begin{table}[H]\centering\footnotesize
\begin{tabularx}{\linewidth}{@{}p{1.5cm} p{3.2cm} X@{}}
\toprule
Day & Goal & Exit criterion \\
\midrule
Thu 17 & We can command the drone, and we know the course & Bench test with propellers off passes on all four axes; piloted measurement flight logged; hover throttle known; course surveyed; overnight training launched \\
Fri 18 & We can see and count gates & Detector working on real gates; gate counter verified offline against recorded footage; autonomous hover; station-keeping in front of one gate \\
Sat 19 & \textbf{A complete lap, by any means} & One clean Plan~B lap. This puts a floor under the score. Plan~A begins with the speed limiter at minimum \\
Sun 20 & Rehearsal and freeze & Both plans rehearsed under run conditions; gates re-measured; configuration frozen; \textbf{Plan~A go/no-go} \\
Mon 21 & \textbf{Scored} & First slot: most reliable configuration, flown slowly, to bank gates. Then step up \\
Tue 22 & \textbf{Scored} & Improve on the banked score. Never risk the drone before a valid run exists \\
\bottomrule
\end{tabularx}
\end{table}

\textbf{Plan~A go/no-go, Sunday.} All three must hold:
\begin{enumerate}
\item it has flown at least three gates autonomously on the real drone, more than once;
\item its observations have been checked against the real camera and the real flight-controller
curves;
\item its failure mode is a safe hold, not a flyaway.
\end{enumerate}

\begin{warnbox}[The scoring rule that should drive every decision]
Ranking is by \textbf{maximum number of gates passed}; time is only the tie-break; and it is the best
run across the final two days. A slow run that completes therefore beats a fast run that ends at gate
four. Bank a score first. Push afterwards. Do not let the last day be the first time we have completed
a run.
\end{warnbox}

%========================================================== QUESTIONS
\section{Questions still open for the organizers}

The specification answered several questions that Revision~1 listed. Answered: the interface is MSP
over UART, with demo software provided; Betaflight CLI access covers rates, PIDs and filters but not
reflashing; camera tilt is adjustable in flight; exposure and gain are controllable; no calibration
matrices are provided; a valid run is two laps with no human pilot intervention; manual override
always supersedes autonomy; practice is in scheduled slots.

Still open:
\begin{enumerate}
\item May we assign an ANGLE mode switch, and change \code{msp\_override\_channels\_mask}, the rate
curve and the throttle curve?
\item May we place tape or markers on the gates to identify them? \textbf{A yes largely removes the
gate-counting problem}, so ask early.
\item Who arms the drone for a scored run? Does arming and enabling override before the start line
count as intervention?
\item Is there a maximum run duration? How many scored attempts per team per day?
\item Are gates re-placed between runs, and to what tolerance?
\item Are gates counted only in order --- does skipping one stop the count? Is a run that crashes
after $N$ gates scored as $N$?
\item Which camera sensor mode is the \SI{75}{\degree} figure quoted for, and can we change modes?
\item Are all four drones on identical firmware and configuration? Is a barometer fitted?
\item Where do we get the camera and serial demo applications?
\item Exact as-built gate heights, and the upper opening height of the double gate.
\end{enumerate}
